\documentclass[11pt,a4paper]{article}

\usepackage[utf8]{inputenc}
\usepackage[T1]{fontenc}
\usepackage[spanish]{babel}
\usepackage{amssymb}
\usepackage[margin=2.4cm]{geometry}
\usepackage{booktabs}
\usepackage{array}
\usepackage{enumitem}
\usepackage{xcolor}
\usepackage{titlesec}
\usepackage{tcolorbox}

\titleformat{\section}{\large\bfseries}{\thesection.}{0.6em}{}
\titlespacing{\section}{0pt}{13pt}{5pt}
\titleformat{\subsection}{\bfseries}{}{0em}{}
\titlespacing{\subsection}{0pt}{9pt}{3pt}
\setlist[itemize]{topsep=3pt,itemsep=2pt,leftmargin=1.4em}
\setlist[enumerate]{topsep=3pt,itemsep=3pt,leftmargin=1.6em}
\newcolumntype{L}[1]{>{\raggedright\arraybackslash}p{#1}}
\newcolumntype{C}[1]{>{\centering\arraybackslash}p{#1}}

\pagestyle{empty}

\usepackage{setspace}
\onehalfspacing

\begin{document}

\begin{center}
{\LARGE\bfseries Proyecto: un clasificador de imágenes propio}\\[4pt]
{\Large IPCC - ML}
\end{center}


\noindent\textbf{Grupos de 4 personas} \hfill
\textbf{Entrega:} un video de 5 minutos \hfill
\textbf{Plazo:} 26 de Octubre 10:00am

\vspace{0.5cm}
\hrule
\vspace{1cm}

\section{El problema}

En clase entrenamos un clasificador de piedra, papel y tijera con un conjunto de imágenes sintéticas. El proyecto consiste en recorrer el
mismo proceso completo, pero con un conjunto de datos construido por ustedes.

\begin{quote}
\textbf{Elijan tres objetos,  construyan el conjunto de imágenes, entrenen el clasificador y analicen en qué se equivoca.}
\end{quote}

\begin{tcolorbox}[colback=red!4,colframe=red!45,boxrule=0.7pt,title=\textbf{Qué se entrega}]
\textbf{Únicamente un video de 5 minutos.} No se entrega el notebook ni el conjunto de imágenes.

El video es la única evidencia del trabajo: lo que no aparezca ahí no se puede calificar.
Los detalles de formato y entrega están más abajo.
\end{tcolorbox}

\vspace{1cm}

\section{Los tres objetos}

Tres objetos físicos, de libre elección, que puedan fotografiar con un celular. Condiciones:

\begin{itemize}
  \item {Distinguibles entre sí: } Si dos de los tres se parecen mucho, el clasificador los va a confundir, y el error va a quedar concentrado en ese par.
  \item {Disponibles en cantidad o en variedad.} 
\end{itemize}

\begin{tcolorbox}[colback=blue!4,colframe=blue!20,boxrule=0.6pt, coltitle=black, title=\textbf{Ejemplos de conjuntos de tres:}]
\begin{itemize}[label*=-]
	\item Lápiz / borrador / sacapuntas
	\item Tres frutas
	\item Llaves / audífonos / cargador
	\item Tres herramientas.
\end{itemize} 
\end{tcolorbox}

\vspace{1cm}

\section{El conjunto de imágenes}

\begin{itemize}
  \item \textbf{60 imágenes por objeto}, 180 en total para \textbf{entrenamiento}. Entre los cuatro son 45 fotos cada uno.
  \item \textbf{Varíen las condiciones deliberadamente:} fondo, iluminación, distancia, ángulo y,
        si decidieron trabajar con categorías, el ejemplar concreto del objeto. Un conjunto de 60
        imágenes casi idénticas aporta mucho menos que 60 variadas.
  \item Organícenlas en tres carpetas con el nombre de cada objeto. Recuerden que la etiqueta de
        cada imagen es el nombre de su carpeta.
  \item \textbf{10 imágenes por objeto} para el test, tomadas en condiciones
        distintas al resto (otro lugar, otro momento, otra luz). Son el conjunto de prueba, y lo van a usar en vivo durante el video.
\end{itemize}

%\begin{tcolorbox}[colback=blue!4,colframe=blue!40,boxrule=0.6pt]
%\textbf{Esta es la parte donde se decide el proyecto.} En clase el modelo falló con la foto de su
%mano por cómo estaba construido el conjunto de entrenamiento, no por la arquitectura de la red.
%Aquí el conjunto lo construyen ustedes, con las decisiones que eso implica.
%\end{tcolorbox}

\vspace{1cm}

\section{Entrenamiento y experimentos}

\begin{enumerate}
	\item Reutilicen el notebook de la clase sustituyendo las imágenes por las suyas. 
	
	\begin{tcolorbox}[colback=blue!4,colframe=blue!20,boxrule=0.6pt, coltitle=black, title=\textbf{Tres clases}]
		Noten que con tres clases, \textbf{responder al azar acierta una de cada tres veces}: un 33\,\% es
		el punto de partida, no un resultado. Apunten a superar el 90\,\%.
	\end{tcolorbox}
	
	\item Después, modifiquen \textbf{una sola cosa}, reentrenen y registren el efecto sobre el accuracy. \textit{Hagan dos modificaciones distintas (no al mismo tiempo). }
	
\end{enumerate}

\vspace{1cm}

\section{El video}

\subsection{Formato y entrega}

\begin{itemize}
  \item \textbf{Duración: 5 minutos.} Es poco tiempo, ensayen antes de grabar.
  \item \textbf{Hablan los cuatro.} Repártanse el contenido por bloques.
  \item \textbf{Se envía a los correos:} mbastidaso@unal.edu.co y jfgarciava@unal.edu.co
  \item \textbf{Nombre del archivo:} los cuatro apellidos en orden alfabético, separados por guion bajo. Por ejemplo: \texttt{Cardona\_Gómez\_Restrepo\_Zapata.mp4}
\end{itemize}

\subsection{Qué debe aparecer}

\begin{enumerate}
  \item \textbf{El conjunto de imágenes.} Abran las carpetas y recorran varias imágenes de cada objeto, de modo que se vea cuántas son y qué tan variadas quedaron. 
  \item \textbf{El código, celda por celda.} Es el componente de mayor peso en la calificación.
  \item \textbf{Los resultados:} el accuracy y la matriz de confusión.
  \item \textbf{La prueba en vivo:} pasen 5 imágenes del conjunto de prueba por el modelo mientras graban, sin cortes, y comenten qué responde en cada una.
  \item \textbf{El análisis:} qué par de objetos confunde, por qué creen que ocurre, y en qué situación usarían o no este clasificador.
\end{enumerate}

\begin{tcolorbox}[colback=black!4,colframe=black!45,boxrule=0.6pt]
\noindent \textbf{Nota:}  Con 5 minutos no alcanza para detenerse en todas las celdas por igual, así que hay que jerarquizar, de las celdas que importan digan tres cosas: \textbf{qué hace}, \textbf{por qué} y \textbf{qué pasaría si se quitara}. 
\end{tcolorbox}

\subsection{Cómo grabar}

\begin{tabular}{L{3.2cm} L{11cm}}
\toprule
Windows & \textbf{Windows + G} abre la barra de grabación. Viene instalada. \\
macOS & \textbf{Cmd + Shift + 5}, seleccionar grabar pantalla. Viene instalado. \\
Alternativa & Abrir una reunión de Meet, compartir pantalla y grabarla. \\
\bottomrule
\end{tabular}

\begin{tcolorbox}[colback=black!4,colframe=black!45,boxrule=0.6pt]
\textbf{Si el modelo falla, muéstrenlo.} No corten la prueba en vivo ni la repitan hasta que
salga bien. Un error mostrado y explicado vale más puntos que una ejecución perfecta sin análisis:
entender por qué un modelo se equivoca es precisamente el objetivo del curso.
\end{tcolorbox}

\section{Rúbrica}

\begin{center}
	\renewcommand{\arraystretch}{1.5}
\begin{tabular}{L{10cm} c}
\toprule
\textbf{Criterio} & \textbf{Pts}  \\
\midrule

\textbf{Explicación del código} & 35 \\

\textbf{Conjunto de imágenes} & 20 \\

\textbf{Prueba en vivo y análisis del error}
& 25 \\

\textbf{Experimentos}
& 10 \\

\textbf{Realización del video}
& 10 \\

\midrule
\textbf{Total} & \textbf{100}  \\
\bottomrule
\end{tabular}
\end{center}

\begin{tcolorbox}[colback=black!4,colframe=black!45,boxrule=0.6pt]
\textbf{El accuracy del modelo no otorga puntos por sí mismo.} Un trabajo con 50\,\% bien explicado califica por encima de uno con 99\,\% donde el código se lee sin entender nada. Lo que se evalúa es el razonamiento sobre lo que hicieron.
\end{tcolorbox}

\newpage 

\section*{Anexo: reducir el tamaño de las imágenes}

Una fotografía de celular pesa alrededor de 3\,MB, de modo que las 180 del proyecto suman más de medio gigabyte. Subir eso a Drive y leerlo desde Colab es lento sin necesidad.

La red entrena con imágenes de 128 píxeles de lado, así que toda la resolución adicional se descarta de todos modos durante el preprocesamiento. Conviene reducirlas antes:

\begin{tcolorbox}[colback=black!3,colframe=black!45,boxrule=0.6pt,fontupper=\ttfamily\small]
from PIL import Image\\
from pathlib import Path\\[4pt]
ORIGINALES = Path("/content/drive/MyDrive/mis\_fotos")\\
REDUCIDAS  = Path("imagenes")\\[4pt]
for foto in ORIGINALES.rglob("*"):\\
\hspace*{1em}if foto.suffix.lower() not in [".jpg", ".jpeg", ".png"]:\\
\hspace*{2em}continue\\
\hspace*{1em}destino = REDUCIDAS / foto.parent.name\\
\hspace*{1em}destino.mkdir(parents=True, exist\_ok=True)\\
\hspace*{1em}img = Image.open(foto).convert('RGB')\\
\hspace*{1em}img.thumbnail((256, 256))\\
\hspace*{1em}img.save(destino / (foto.stem + ".jpg"), quality=85)\\[4pt]
print("Listo. Usar CARPETA = Path('imagenes')")
\end{tcolorbox}

\vspace{4pt}
\noindent\emph{Una comprobación importante para ustedes:} abrir una imagen original y su versión reducida lado a lado. ¿Son visualmente equivalentes?

\end{document}
